Software developers and malware authors widely use code obfuscation to conceal program logic.
In real-world scenarios, these actors rarely rely on isolated techniques. Instead, they stack multiple
transformations together to build resilient, layered defenses. While large language models (LLMs)
can effectively reason about clean source code, their performance drops significantly when they
encounter these complex defenses. Despite stacked transformations being the realistic standard
in practice, current literature almost exclusively evaluates models against individual, isolated
obfuscations. Existing approaches fundamentally fail to assess and effectively handle layered
techniques, and learning to counter individual transformations does not automatically help models
handle stacked ones. To address this critical gap, we evaluate model merging as a solution for realistic stacked
obfuscation. We compare this approach against three fine-tuning baselines, which achieve strong
performance on familiar transformations but struggle to generalize beyond their training distribution
and can exhibit catastrophic forgetting of untrained reasoning capabilities. Model merging mitigates
these limitations and achieves the strongest overall balance between performance on stacked
obfuscations and preservation of untrained reasoning capabilities. Merged models significantly outperform their base models on obfuscated code and, surprisingly,
also improve prediction on the original, clean programs. They additionally preserve performance on
an untrained reverse-reasoning task, where the fine-tuning baselines frequently degrade. By
successfully combining specific adaptations with broader reasoning capabilities, these results
demonstrate that merging independently learned models is a highly effective strategy for analyzing
complex, realistically obfuscated code.